\section{评估、治理与部署}

\subsection{多层评估矩阵}

一个完整的评估流程应包含自动 benchmark、trajectory review 和 stress test：

\begin{itemize}
    \item `score`、`TD error`、`episode length` 等数值指标
    \item 手动审阅 failure trajectory / reward hack
    \item 噪声注入、belief shift、环境 perturbation 测试稳定性
\end{itemize}

\begin{table}[H]
\centering
\begin{tabular}{p{0.25\textwidth} p{0.25\textwidth} p{0.2\textwidth}}
\toprule
层级 & 指标 & 触发策略 \\
\midrule
Training eval & Q target, loss, entropy & Divergence or reward plateau \\
Safety audit & reward hack detection, adversarial action & anomalous trajectories \\
Deployment sim & latency, repeatability, replicates & hardware change / rollout \\
\bottomrule
\end{tabular}
\caption{评估与治理矩阵}
\end{table}

\begin{knowledgebox}{日志监控建议}
记录 TD error distribution、max Q target、exploration ratio；设置 \texttt{td\_error > threshold} 触发 replay augmentation，\texttt{reward variance > 2σ} 触发 reward shaping review。
\end{knowledgebox}

\subsection{部署与治理 checklist}

\begin{warningbox}{上线前的三道“门”}
\begin{enumerate}
    \item Replay buffer 只能包含经过 human-in-the-loop 校验的 transitions。
    \item Checkpoint 避免 reward hack，如 reward clipping + sanity check trajectories。
    \item 设定 clear rollback 条件（loss divergence、reward drop > threshold）。
\end{enumerate}
\end{warningbox}

\begin{figure}[H]
\centering
\includegraphics[width=0.86\textwidth,page=2]{06_cs224r_qlearning_2025.pdf}
\caption{Slide：Q-Learning deployment guardrails 与 monitoring stack}
\end{figure}

\subsection{Evidence Matrix 与操作时间线}
\subsection{Evidence Matrix 与操作时间线}

为每轮训练记录 evidence matrix，可帮助快速定位 failure root cause。同一矩阵也能映射到 action timeline：

\begin{table}[H]
\centering
\begin{tabular}{p{0.28\textwidth} p{0.58\textwidth}}
\toprule
Evidence & Action \\
\midrule
Prediction explosion & 回滚至低 learning rate checkpoint + 增加 target damping \\
Reward spike & 检查 reward shaping logs，暂停 training 并审查 newly added shaping term \\
TD error drift & 触发 replay augmentation、增加 PER alpha \\
Policy collapse (success rate drop) & Enforce early rollback + high-priority fail buffer replay \\
\bottomrule
\end{tabular}
\caption{Evidence matrix 与对应行动时间线}
\end{table}

\begin{knowledgebox}{Actions should be traceable}
为每次 triggered action 记录 context（checkpoint id、evidence source、actor responsible），方便审计和 replay review。
\end{knowledgebox}

\subsection{本章小结}

评估矩阵和治理 checklist 保证 Q-Learning 算法不会在部署后被 reward hack 或 value explosion 破坏。

